\chapter{Montaje Experimental}
\thispagestyle{fancy}
\label{ch:montaje}

\subsection{Objetivos}

\begin{itemize}
    \item Realizar mediciones directas de masa utilizando la balanza granataria, y de dimensiones lineales (longitud y diámetro) empleando el vernier y el micrómetro. 
    \item Determinar la densidad de diversos sólidos regulares a partir de las magnitudes medidas, utilizando modelos geométricos de volumen. 
    \item Aplicar el método de propagación de incertidumbres bajo el criterio del peor escenario para evaluar la precisión de las magnitudes indirectas. 
    \item Evaluar las fuentes de error sistemático y aleatorio en el proceso experimental, identificando la contribución más significativa a la incertidumbre total. 
\end{itemize}

\subsection{Montaje}

Para realizar esta experiencia, se midió la masa, longitud y diámetro de cinco sólidos: una moneda, un cilindro plateado, un cilindro dorado, un corcho y un cilindro hueco. Para estas mediciones, se hizo uso de los siguientes instrumentos:
 
\begin{itemize}
    \item \textbf{Balanza granataria}: con apreciación de $0,01 \text{ g}$. 
    \item \textbf{Vernier}: con apreciación de $0,005 \text{ cm}$. 
    \item \textbf{Micrómetro}: con apreciación de $0,001 \text{ cm}$ ($0,01 \text{ mm}$).
\end{itemize}

Para las medidas con el vernier, se realizaron entre 10 y 12 repeticiones de longitud y diámetro por cada sólido para obtener un valor promedio representativo y minimizar el impacto de errores aleatorios.  En el caso del micrómetro, se tomaron 12 mediciones de espesor en la moneda, controlando cuidadosamente la presión del tornillo para evitar errores sistemáticos por deformación. 

Para optimizar la técnica y asegurar la repetibilidad, se utilizó el siguiente sistema para cada sólido:
\begin{itemize}
    \item \textbf{Cilindros plateado y dorado}: Se midió la longitud y diámetro con el vernier rotando el sólido y variando los puntos de contacto, asegurando siempre la perpendicularidad del instrumento respecto al eje del cilindro. 
    \item \textbf{Cilindro hueco}: Se midió la longitud y el diámetro externo con el vernier. Para el diámetro interno, se emplearon las mordazas superiores del instrumento, garantizando un ajuste firme contra las paredes internas.
    \item \textbf{Corcho}: Se midió con extremo cuidado para evitar deformaciones por presión. Debido a su baja masa, esta se determinó por el método de diferencia: se registró la masa de la moneda sola y luego la del conjunto moneda-corcho, obteniendo la masa neta por sustracción. 
    \item \textbf{Moneda}: Se utilizó el micrómetro para medir el espesor en distintos puntos de la superficie para reducir errores aleatorios. Para el diámetro, dada su magnitud, se utilizó el vernier. 
\end{itemize}
